\section{Closed-Loop Policy Evaluation (RQ2, part 1)}
\label{sec:closed-loop}

\begin{sloppypar}
Section~\ref{sec:characterization} characterizes the offline target.
This section reports the actual closed-loop cache-policy behavior it will
be compared against in Section~\ref{sec:linkage} (Table~\ref{tab:tier1-closed-loop}). Source: Tier-1
production closed-loop evaluation (repository path
\texttt{analysis/closed\_loop\_tier1\_evidence\_20260914/}; RUN\_ID
\texttt{20260914T023445Z\_8a4cd32402a1}; 14/14 validity gates
passed; 230/230 planned executions complete, 0 failed). Scope: 5 families
$\times$ 2 capacities (32, 128) $\times$ 4 policies (LRU, MRU, random-mean
over 20 seeds, SIEVE~\cite{zhang2024sieve}), per
Section~\ref{sec:experimental-methodology}.
\end{sloppypar}

\begin{table}[t]
  \centering
  \footnotesize
  \caption{Tier-1 closed-loop production evaluation: LRU miss ratio and relative miss-ratio difference vs.\ LRU for MRU, mean-random (20 seeds), and SIEVE. Source: \texttt{analysis/closed\_loop\_tier1\_evidence\_20260914/} (230/230 complete, 14/14 gates passed). MetaCDN is validation-window evidence; all other families are test-window. Negative values mean the policy beats LRU.}
  \label{tab:tier1-closed-loop}
  \begin{tabularx}{\columnwidth}{@{}llrrrrr@{}}
    \toprule
    family & cap. & LRU miss ratio & MRU $\Delta$ & random $\Delta$ & SIEVE $\Delta$ & random std \\
    \midrule
    metacdn      & 32  & 0.5792 & +47.37\% & +5.38\% & +0.81\%  & 0.0013 \\
    metacdn      & 128 & 0.5490 & +54.93\% & +2.98\% & +0.36\%  & 0.0014 \\
    twemcache    & 32  & 0.7511 & +27.12\% & +4.19\% & +8.47\%  & 0.0025 \\
    twemcache    & 128 & 0.6205 & +52.98\% & +7.78\% & +15.29\% & 0.0023 \\
    metakv       & 32  & 0.7410 & +10.02\% & +0.47\% & +0.03\%  & 0.0011 \\
    metakv       & 128 & 0.7290 & +8.57\%  & $-$2.07\% & $-$4.22\% & 0.0015 \\
    cloudphysics & 32  & 0.9836 & +1.65\%  & +0.28\% & +0.20\%  & 0.0010 \\
    cloudphysics & 128 & 0.9561 & +4.48\%  & +0.85\% & +1.44\%  & 0.0014 \\
    wiki2018     & 32  & 1.0000 & 0.00\%   & 0.00\%  & 0.00\%   & 0.0000 \\
    wiki2018     & 128 & 1.0000 & 0.00\%   & 0.00\%  & 0.00\%   & 0.0000 \\
    \bottomrule
  \end{tabularx}
\end{table}


\textbf{LRU is the best policy in 9 of 10 cells.} The sole exception is
metakv/cap128, where SIEVE (miss ratio 0.6982) beats LRU (0.7290) by
4.22\% relative -- the only cell in the entire evaluation where a
non-LRU policy wins outright. \textbf{MRU is the worst policy in every
non-degenerate cell}, by a wide margin (relative miss-ratio increase over
LRU ranging from 1.65\% at alibaba-block/cap32, a near-zero-signal regime,
up to 54.93\% at metacdn/cap128). \textbf{Random-seed variability is small
relative to policy gaps} throughout: the largest random-policy standard
deviation observed is 0.0025 (twemcache/cap32), an order of magnitude
below the smallest genuine policy gap reported in the table.

\textbf{wiki2018 is a complete, exact tie across every policy}: miss ratio
1.0 for LRU, MRU, SIEVE, and every one of random's 20 seeds, at both
capacities, with random's standard deviation exactly 0.0. This is not
approximate degeneracy -- it is byte-for-byte identical across all 4
policies and 20 seeds. Section~\ref{sec:mechanisms} shows this follows
deductively from zero reuse events in the scored window, not as a
statistical artifact of this particular closed-loop run.

\textbf{alibaba-block is a near-tie regime at both capacities}: every
policy's relative gap to LRU stays under 4.5\%, an order of magnitude
smaller than metacdn's or twemcache's gaps. \textbf{The metakv/cap128 cell is the one genuine exception to "LRU
dominates"}, and there both SIEVE (miss ratio 0.6982) and random-mean
(approximately 0.7137, relative gap $-2.07\%$) beat LRU (0.7290), while
MRU remains worst; Section~\ref{sec:mechanisms} investigates, but does not
fully resolve, why. metacdn's results should be read as validation-window
evidence throughout (Section~\ref{sec:experimental-methodology});
twemcache's are test-window.

No universal middle-of-the-pack ranking is claimed for SIEVE vs.\ random:
SIEVE beats random-mean at metacdn and metakv/cap32, but random-mean beats
SIEVE at twemcache (both capacities), alibaba-block (both capacities), and
metakv/cap128 -- the relative order of these two policies is
family-dependent, not fixed. Only one ordering statement holds across
every non-degenerate cell: \textbf{MRU is worst everywhere}. Among this
four-policy comparator set, LRU is best in 9 of 10 cells and falls only as
low as third place, at metakv/cap128, in the tenth.

\subsection{Robustness to an Expanded Comparator Set}
\label{sec:expanded-comparators}

\begin{sloppypar}
LRU/MRU/random/SIEVE is a deliberately small, diagnostic comparator set,
not a full policy-landscape comparison (Section~\ref{sec:related-work}
surveys the broader landscape of classical and modern replacement
policies \lafc\ interfaces with). To test whether the conclusions above
are an artifact of that narrow set, we add three representative stronger
baselines under the identical Tier-1 protocol (same 5 families, same
capacities 32/128, same scored windows and trace hashes): \textbf{ARC}
\cite{megiddo2003arc} (adaptive recency/frequency), \textbf{LIRS}
\cite{jiang2002lirs} (inter-reference-recency), and \textbf{S3-FIFO}
\cite{yang2023s3fifo} (modern lightweight FIFO-family, using its own
published default ratios, not tuned to any trace). All three are
deterministic, reference C implementations from libCacheSim, the same
simulator family as the SIEVE reference implementation already used in
Table~\ref{tab:tier1-closed-loop}, independently validated before
production use (hand-derived tiny-trace tracing for ARC, a cross-check of
the library's own LRU engine against this manuscript's already-trusted
LRU implementation, and a scan-resistance property check). A fourth
candidate, W-TinyLFU, was pre-registered but dropped after validation
found its reference implementation non-functional at capacity 32 under
its own published default configuration; see this experiment's
\texttt{COMPARATOR\_PROTOCOL.md} for the full pre-registration and this
finding.
\end{sloppypar}

\input{tables/table_problem5_expanded_closed_loop}

\begin{figure}[t]
  \centering
  \includegraphics[width=\columnwidth]{figures/figure_problem5_policy_spread_vs_discriminativeness.pdf}
  \caption{Closed-loop policy spread (max$-$min miss ratio across all 7
  policies) versus offline discriminative fraction at $H=16$, one point per
  family$\times$capacity cell. The offline-discriminativeness$\leftrightarrow$
  closed-loop-separation association (Section~\ref{sec:linkage}) is
  essentially unchanged when the policy landscape is expanded ($r=0.96$
  with the original 4-policy spread vs.\ $r=0.96$ with the expanded
  7-policy spread, $n=10$).}
  \label{fig:policy-spread}
\end{figure}

Under the expanded 7-policy set, LRU is strictly the best policy in 4 of
10 cells (down from 9 of 10 among the original four). However, four of
the six cells where a new policy edges out LRU do so by less than 3.2\%
relative miss ratio -- comparable in size to gaps this manuscript already
treats as near-ties elsewhere (e.g.\ alibaba-block's sub-4.5\% policy
gaps) -- so most of this apparent reordering reflects fine-grained
jockeying among already-close policies, not a qualitative reversal. The
one cell with a clearly large margin is the already-known exception,
\textbf{metakv/cap128}, where the gap to LRU widens under the broader set
(S3-FIFO's miss ratio is 8.07\% below LRU's, versus SIEVE's previously
reported 4.22\%); no new family/capacity cell crosses that magnitude. LRU
remains worst-case-safe and near-best almost everywhere; it is not
displaced as the practical default, but it is no longer accurate to say it
wins outright in nearly every cell once the comparator set is broadened.
Most importantly for RQ3, the link between offline discriminativeness and
closed-loop policy separation is not an artifact of the narrow original
comparator set: it is essentially identical whether policy separation is
measured across 4 or 7 policies (Figure~\ref{fig:policy-spread}).
